\section{Perception and Spatial Reasoning}
\label{sec:perception}

\subsection{Sensing the World}
\begin{itemize}
  \item NPCs need to ``see'' and ``hear'' to react believably; usually this
        is simulated with vision cones, line-of-sight checks and sound
        radii~\cite{millington2009}.
  \item Designers deliberately limit perception: an NPC that always knows
        where the player is feels unfair and breaks stealth
        gameplay~\cite{yannakakis2018}.
  \item Stealth games make perception visible to the player (alert
        indicators), turning AI awareness into a game mechanic.
  \item Memory matters: NPCs remember the last known position of the player
        and search there instead of instantly forgetting~\cite{isla2001}.
\end{itemize}

\subsection{Pathfinding}
\begin{itemize}
  \item Finding a route from A to B is the most common AI task in games.
  \item The A* algorithm~\cite{hart1968} remains the standard: it explores
        promising directions first by estimating the remaining distance.
  \item Explain intuitively with a city-map analogy; no formulas needed.
  \item \textbf{[Figure idea]} Grid with obstacles showing the explored
        area and the final path.
  \item Worlds are simplified for the search: waypoint graphs or
        \emph{navigation meshes} (walkable areas made of polygons) are now
        the industry norm~\cite{millington2009}.
  \item Large worlds use hierarchical pathfinding: plan between regions
        first, then refine locally.
\end{itemize}

\subsection{Steering and Movement}
\begin{itemize}
  \item Pathfinding gives a route; steering decides how to move along it
        smoothly and avoid other characters~\cite{reynolds1999}.
  \item Reynolds' ``boids'' show that three simple rules (separation,
        alignment, cohesion) produce realistic flocks and
        crowds~\cite{reynolds1987}.
  \item Used for crowds, animal herds and squads moving in formation.
\end{itemize}

\subsection{Tactical Reasoning}
\begin{itemize}
  \item NPCs evaluate positions, not just routes: Where is good cover?
        Where is the enemy likely to come from?
  \item \emph{Influence maps} spread values such as danger or control
        across the map, helping strategy AI decide where to attack or
        defend~\cite{millington2009}.
  \item Tactical position scoring is closely related to utility AI
        (Section~\ref{sec:modern}).
\end{itemize}

\subsection{Summary}
\begin{itemize}
  \item Good perception and movement are often what players notice first;
        clever decisions are wasted if the NPC gets stuck on a door.
  \item These systems are mature and widely shared across games, unlike
        decision making, which varies more by genre.
\end{itemize}
